\documentclass[a4paper]{article}

\usepackage{amsmath,amssymb}
\usepackage{xcolor}
\usepackage{hyperref}
\usepackage{geometry}

\geometry{margin=2.5cm}

\hypersetup{
    colorlinks=true,
    linkcolor=blue,
    urlcolor=blue,
    citecolor=blue
}

\newcommand{\FMM}{\operatorname{FMM}}
\newcommand{\Transformer}{\operatorname{Transformer}}
\newcommand{\Decoder}{\operatorname{Decoder}}
\newcommand{\Encoder}{\operatorname{Encoder}}
\newcommand{\Project}{\operatorname{Project}}
\newcommand{\SensorDecoder}{\operatorname{SensorDecoder}}

\title{Reversible Vector-Valued FMM Wave Fusion for Multiple Sensors}
\date{}
\author{Mohammad Rahmani}

\begin{document}
    \maketitle
    
    
    \section{Single FMM Wave}
        
        A single FMM wave can be written as:
        
        \begin{equation}
            W(s;M,A,\alpha,\beta,\omega)
            =
            M
            +
            A
            \cos
            \biggl(
            \beta
            +
            2
            \arctan
            \biggl(
            \omega
            \tan
            \biggl(
            \frac{s-\alpha}{2}
            \biggr)
            \biggr)
            \biggr)
        \end{equation}
        
        where \(s \in [0,2\pi]\) is the normalized temporal variable of the wave.
        
        \begin{itemize}
            \item \(W(s)\) is the value of the FMM wave at normalized time \(s\).
            
            \item \(M\) is the baseline or vertical offset.
            
            \item \(A\) is the amplitude parameter.
            
            \item \(\alpha\) is the location parameter.
            
            \item \(\beta\) is a phase-shape parameter.
            
            \item \(\omega\) is a shape parameter related to peakedness, concentration, smoothness, and local temporal deformation.
        \end{itemize}
        
        The parameter \(\omega\) is not the wavelength of the FMM wave.
        The horizontal comparability of several FMM waves is obtained by evaluating them on the same normalized domain \(s \in [0,2\pi]\), not by forcing them to have the same \(\omega\).
        
        If the original time variable is \(t' \in [t_0,t_0+T]\), it can be normalized as:
        
        \begin{equation}
            s
            =
            \frac{2\pi(t'-t_0)}{T}
        \end{equation}
    
    
    \section{Sensor-to-FMM Mapping}
        
        Assume that there are \(n\) sensors:
        
        \begin{equation}
            S_t^1,
            S_t^2,
            \ldots,
            S_t^n
        \end{equation}
        
        Each sensor reading is encoded by a sensor-specific VAE encoder into one FMM parameter vector:
        
        \begin{equation}
            \theta_t^i
            =
            \Encoder_i(S_t^i)
            =
            \bigl(
            M_t^i,
            A_t^i,
            \alpha_t^i,
            \beta_t^i,
            \omega_t^i
            \bigr),
            \qquad
            i=1,\ldots,n
        \end{equation}
        
        Each FMM parameter vector generates one sensor-specific wave:
        
        \begin{equation}
            w_t^i(s)
            =
            \FMM
            \bigl(
            s;
            M_t^i,
            A_t^i,
            \alpha_t^i,
            \beta_t^i,
            \omega_t^i
            \bigr)
        \end{equation}
        
        Each sensor may have its own bounded FMM parameter region:
        
        \begin{equation}
            \theta_t^i
            \in
            \mathcal{B}_i
        \end{equation}
        
        where \(\mathcal{B}_i\) is the admissible parameter region for sensor \(i\).
    
    
    \section{Orthonormal Sensor Axes}
        
        To make the composite wave reversible, each sensor is assigned to its own orthonormal axis.
        Let:
        
        \begin{equation}
            \mathbf{e}_1,
            \mathbf{e}_2,
            \ldots,
            \mathbf{e}_n
            \in
            \mathbb{R}^{n}
        \end{equation}
        
        be the sensor basis vectors, with:
        
        \begin{equation}
            \mathbf{e}_i^{\top}
            \mathbf{e}_j
            =
            \delta_{ij}
        \end{equation}
        
        where:
        
        \begin{equation}
            \delta_{ij}
            =
            \begin{cases}
                1,
                &
                i=j
                \\
                0,
                &
                i \neq j
            \end{cases}
        \end{equation}
        
        For example, with three sensors:
        
        \begin{equation}
            \mathbf{e}_{GPS}
            \perp
            \mathbf{e}_{LiDAR}
            \perp
            \mathbf{e}_{IMU}
        \end{equation}
        
        This is the key difference from scalar fusion.
        The sensor waves are not collapsed into one scalar product.
        They are stored as orthogonal components of one vector-valued wave.
    
    
    \section{Reversible Vector-Valued Composite Wave}
        
        The reversible composite wave is defined as:
        
        \begin{equation}
            \mathbf{U}_t(s)
            =
            \sum_{i=1}^{n}
            w_t^i(s)
            \mathbf{e}_i
        \end{equation}
        
        Equivalently:
        
        \begin{equation}
            \mathbf{U}_t(s)
            =
            w_t^1(s)\mathbf{e}_1
            +
            w_t^2(s)\mathbf{e}_2
            +
            \cdots
            +
            w_t^n(s)\mathbf{e}_n
        \end{equation}
        
        In coordinate form:
        
        \begin{equation}
            \mathbf{U}_t(s)
            =
            \begin{bmatrix}
                w_t^1(s)
                \\
                w_t^2(s)
                \\
                \vdots
                \\
                w_t^n(s)
            \end{bmatrix}
            \in
            \mathbb{R}^{n}
        \end{equation}
        
        Therefore, \(\mathbf{U}_t(s)\) is not a scalar wave.
        It is a vector-valued wave whose components preserve sensor identity.
    
    
    \section{Exact Recovery of Sensor Waves}
        
        Because the sensor axes are orthonormal, each sensor-specific wave can be recovered by projection:
        
        \begin{equation}
            w_t^i(s)
            =
            \mathbf{U}_t(s)^{\top}
            \mathbf{e}_i
        \end{equation}
        
        Therefore:
        
        \begin{equation}
            \Project_i
            \bigl(
            \mathbf{U}_t(s)
            \bigr)
            =
            w_t^i(s)
        \end{equation}
        
        For all sensors:
        
        \begin{equation}
            \mathbf{U}_t(s)
            \longrightarrow
            \bigl(
            w_t^1(s),
            w_t^2(s),
            \ldots,
            w_t^n(s)
            \bigr)
        \end{equation}
        
        This recovery is exact at the wave-component level, provided that \(\mathbf{U}_t(s)\) is preserved as a vector-valued wave.
    
    
    \section{Discrete Implementation}
        
        In implementation, the continuous variable \(s\) is sampled on a fixed grid:
        
        \begin{equation}
            \mathcal{G}
            =
            \{
            s_1,
            s_2,
            \ldots,
            s_K
            \}
            \subset
            [0,2\pi]
        \end{equation}
        
        For each time \(t\), the vector-valued wave is represented as:
        
        \begin{equation}
            \mathbf{U}_t
            =
            \begin{bmatrix}
                \mathbf{U}_t(s_1)^{\top}
                \\
                \mathbf{U}_t(s_2)^{\top}
                \\
                \vdots
                \\
                \mathbf{U}_t(s_K)^{\top}
            \end{bmatrix}
            \in
            \mathbb{R}^{K \times n}
        \end{equation}
        
        Each column corresponds to one sensor wave.
        Thus, the sensor identity is preserved during fusion.
    
    
    \section{Optional Interaction Features}
        
        Interaction features may be added, but they must not replace the reversible vector-valued wave.
        
        The average wave may be defined as:
        
        \begin{equation}
            \bar{w}_t(s)
            =
            \frac{1}{n}
            \sum_{i=1}^{n}
            w_t^i(s)
        \end{equation}
        
        For \(n \geq 2\), the pairwise interaction wave may be defined as:
        
        \begin{equation}
            p_t(s)
            =
            \frac{2}{n(n-1)}
            \sum_{i<j}
            w_t^i(s)
            w_t^j(s)
        \end{equation}
        
        An augmented representation can then be written as:
        
        \begin{equation}
            \mathbf{X}_t(s)
            =
            \begin{bmatrix}
                \mathbf{U}_t(s)
                \\
                \bar{w}_t(s)
                \\
                p_t(s)
            \end{bmatrix}
            \in
            \mathbb{R}^{n+2}
        \end{equation}
        
        The first \(n\) components of \(\mathbf{X}_t(s)\) are reversible sensor-specific components.
        The last two components are auxiliary interaction features.
        
        Therefore:
        
        \begin{equation}
            \mathbf{X}_t(s)
            \longrightarrow
            \mathbf{U}_t(s)
            \longrightarrow
            \bigl(
            w_t^1(s),
            w_t^2(s),
            \ldots,
            w_t^n(s)
            \bigr)
        \end{equation}
        
        The auxiliary interaction features may help prediction, but they are not used as the only source for reconstruction.
    
    
    \section{Temporal Prediction}
        
        The Transformer receives a temporal sequence of vector-valued waves:
        
        \begin{equation}
            \mathbf{U}_{t-T},
            \mathbf{U}_{t-T+1},
            \ldots,
            \mathbf{U}_{t}
        \end{equation}
        
        and predicts the next vector-valued wave:
        
        \begin{equation}
            \Transformer
            \bigl(
            \mathbf{U}_{t-T},
            \mathbf{U}_{t-T+1},
            \ldots,
            \mathbf{U}_{t}
            \bigr)
            =
            \hat{\mathbf{U}}_{t+1}
        \end{equation}
        
        If the augmented representation is used, then:
        
        \begin{equation}
            \Transformer
            \bigl(
            \mathbf{X}_{t-T},
            \mathbf{X}_{t-T+1},
            \ldots,
            \mathbf{X}_{t}
            \bigr)
            =
            \hat{\mathbf{X}}_{t+1}
        \end{equation}
        
        and the predicted vector-valued wave is extracted from the first \(n\) components of \(\hat{\mathbf{X}}_{t+1}\).
    
    
    \section{Recovery After Forecasting}
        
        After forecasting, each sensor-specific predicted wave is recovered by projection:
        
        \begin{equation}
            \hat{w}_{t+1}^i(s)
            =
            \hat{\mathbf{U}}_{t+1}(s)^{\top}
            \mathbf{e}_i
        \end{equation}
        
        For all sensors:
        
        \begin{equation}
            \hat{\mathbf{U}}_{t+1}(s)
            \longrightarrow
            \bigl(
            \hat{w}_{t+1}^1(s),
            \hat{w}_{t+1}^2(s),
            \ldots,
            \hat{w}_{t+1}^n(s)
            \bigr)
        \end{equation}
        
        Then each predicted wave is mapped to its FMM parameter vector:
        
        \begin{equation}
            \hat{w}_{t+1}^i(s)
            \longrightarrow
            \hat{\theta}_{t+1}^i
        \end{equation}
        
        where:
        
        \begin{equation}
            \hat{\theta}_{t+1}^i
            =
            \bigl(
            \hat{M}_{t+1}^i,
            \hat{A}_{t+1}^i,
            \hat{\alpha}_{t+1}^i,
            \hat{\beta}_{t+1}^i,
            \hat{\omega}_{t+1}^i
            \bigr)
        \end{equation}
        
        Finally, the sensor-specific VAE decoder reconstructs the sensor reading:
        
        \begin{equation}
            \hat{S}_{t+1}^i
            =
            \SensorDecoder_i
            \bigl(
            \hat{\theta}_{t+1}^i
            \bigr)
        \end{equation}
        
        Thus:
        
        \begin{equation}
            \hat{\mathbf{U}}_{t+1}
            \longrightarrow
            \bigl(
            \hat{S}_{t+1}^1,
            \hat{S}_{t+1}^2,
            \ldots,
            \hat{S}_{t+1}^n
            \bigr)
        \end{equation}
    
    
    \section{Training Losses}
        
        The vector-valued wave prediction loss is:
        
        \begin{equation}
            \mathcal{L}_{U}
            =
            \frac{1}{K}
            \sum_{k=1}^{K}
            \left\lVert
                \mathbf{U}_{t+1}(s_k)
            -
                \hat{\mathbf{U}}_{t+1}(s_k)
            \right\rVert_2^2
        \end{equation}
        
        The sensor-specific wave reconstruction loss is:
        
        \begin{equation}
            \mathcal{L}_{w}
            =
            \sum_{i=1}^{n}
            \frac{1}{K}
            \sum_{k=1}^{K}
            \bigl(
            w_{t+1}^i(s_k)
            -
            \hat{w}_{t+1}^i(s_k)
            \bigr)^2
        \end{equation}
        
        The sensor reconstruction loss is:
        
        \begin{equation}
            \mathcal{L}_{S}
            =
            \sum_{i=1}^{n}
            \left\lVert
                S_{t+1}^i
            -
                \hat{S}_{t+1}^i
            \right\rVert_2^2
        \end{equation}
        
        Because \(\alpha\) and \(\beta\) are angular parameters, their errors should be measured with a circular loss:
        
        \begin{equation}
            d_{\mathbb{S}^1}(x,y)
            =
            1
            -
            \cos(x-y)
        \end{equation}
        
        For the non-angular parameters, define:
        
        \begin{equation}
            r_t^i
            =
            \bigl(
            M_t^i,
            A_t^i,
            \omega_t^i
            \bigr)
        \end{equation}
        
        and:
        
        \begin{equation}
            \hat{r}_{t+1}^i
            =
            \bigl(
            \hat{M}_{t+1}^i,
            \hat{A}_{t+1}^i,
            \hat{\omega}_{t+1}^i
            \bigr)
        \end{equation}
        
        The FMM parameter prediction loss is:
        
        \begin{equation}
            \mathcal{L}_{\theta}
            =
            \sum_{i=1}^{n}
            \biggl(
            \left\lVert
                r_{t+1}^i
            -
                \hat{r}_{t+1}^i
            \right\rVert_2^2
            +
            \lambda_{\alpha}
            d_{\mathbb{S}^1}
            \bigl(
            \alpha_{t+1}^i,
            \hat{\alpha}_{t+1}^i
            )
            +
            \lambda_{\beta}
            d_{\mathbb{S}^1}
            \bigl(
            \beta_{t+1}^i,
            \hat{\beta}_{t+1}^i
            )
            \biggr)
        \end{equation}
        
        The total loss is:
        
        \begin{equation}
            \mathcal{L}
            =
            \lambda_U
            \mathcal{L}_{U}
            +
            \lambda_w
            \mathcal{L}_{w}
            +
            \lambda_S
            \mathcal{L}_{S}
            +
            \lambda_{\theta}
            \mathcal{L}_{\theta}
        \end{equation}
    
    
    \section{Novelty Detection}
        
        Novelty detection can be performed in three related spaces:
        
        \begin{enumerate}
            \item The vector-valued wave space.
            
            \item The FMM parameter space.
            
            \item The raw sensor space.
        \end{enumerate}
        
        The wave-space novelty score for sensor \(i\) is:
        
        \begin{equation}
            \eta_{w,t+1}^i
            =
            \sqrt{
                \frac{1}{K}
                \sum_{k=1}^{K}
                \bigl(
                w_{t+1}^i(s_k)
                -
                \hat{w}_{t+1}^i(s_k)
                \bigr)^2
            }
        \end{equation}
        
        The FMM-parameter novelty score for sensor \(i\) is:
        
        \begin{equation}
            \eta_{\theta,t+1}^i
            =
            \left\lVert
                r_{t+1}^i
            -
                \hat{r}_{t+1}^i
            \right\rVert_2
            +
            d_{\mathbb{S}^1}
            \bigl(
            \alpha_{t+1}^i,
            \hat{\alpha}_{t+1}^i
            )
            +
            d_{\mathbb{S}^1}
            \bigl(
            \beta_{t+1}^i,
            \hat{\beta}_{t+1}^i
            )
        \end{equation}
        
        The raw sensor novelty score for sensor \(i\) is:
        
        \begin{equation}
            \rho_{t+1}^i
            =
            \left\lVert
                S_{t+1}^i
            -
                \hat{S}_{t+1}^i
            \right\rVert_2
        \end{equation}
    
    
    \section{Important Modeling Note}
        
        The reversible composite wave is:
        
        \begin{equation}
            \mathbf{U}_t(s)
            =
            \sum_{i=1}^{n}
            w_t^i(s)
            \mathbf{e}_i
        \end{equation}
        
        This representation is reversible at the sensor-wave level because the sensor axes are orthonormal.
        Each sensor wave is recovered by projection:
        
        \begin{equation}
            w_t^i(s)
            =
            \mathbf{U}_t(s)^{\top}
            \mathbf{e}_i
        \end{equation}
        
        In contrast, a scalar composite wave such as:
        
        \begin{equation}
            u_t(s)
            =
            a
            \frac{1}{n}
            \sum_{i=1}^{n}
            w_t^i(s)
            +
            b
            \frac{2}{n(n-1)}
            \sum_{i<j}
            w_t^i(s)
            w_t^j(s)
        \end{equation}
        
        is not generally reversible.
        It can be useful as an auxiliary feature, but it should not replace the vector-valued composite wave if sensor-specific recovery is required.
    
    
    \section{Summary}
        
        The proposed method maps each sensor reading into an FMM parameter vector using a sensor-specific VAE encoder.
        Each parameter vector generates a sensor-specific FMM wave.
        Instead of collapsing all waves into one scalar wave, the method places each sensor wave on an orthonormal sensor axis.
        The resulting composite wave is a vector-valued wave.
        Sensor-specific waves are recoverable by projection from the vector-valued composite wave.
        A Transformer predicts the future vector-valued wave, and each predicted sensor wave is then projected, mapped to FMM parameters, and decoded back into the corresponding sensor reading.
        
        For Python implementation, see:
        \url{https://github.com/FMMGroupVa/fmmpy}
        
        For the R package, see:
        \url{https://cran.r-project.org/web/packages/FMM/vignettes/FMMVignette.html}

\end{document}